\documentclass[12pt]{article}
\usepackage[margin=1in]{geometry}
\usepackage{amsmath, amssymb}
\usepackage{hyperref}
\usepackage{parskip}
\usepackage{microtype}
\usepackage{lmodern}
\usepackage{booktabs}

\hypersetup{colorlinks=true, linkcolor=black, urlcolor=blue, citecolor=blue}

\title{Hybrid Ground Truth Realism \\ Pass 1 --- Conceptual Foundation}
\author{Dandelion Engineering}
\date{\today}

\begin{document}
\maketitle
\tableofcontents
\newpage

\section{Introduction}

This project asks whether a measuring instrument the field of neuroscience relies on is
actually measuring what everyone assumes it measures. That is the whole of it, and the
rest of this document exists so that you can hold that sentence with real content behind
each word.

The instrument is a \emph{benchmark}. Neuroscientists have an algorithmic step they cannot
easily check --- deciding which electrical spike in a recording came from which neuron ---
and they check it by building recordings where they secretly know the answer. They take a
real recording and inject synthetic spikes into it at times of their own choosing. Then
they run the algorithm and count how many of the planted spikes it found. The planted
spikes are the answer key.

The maintainers who built the field's standard version of that answer key wrote, in print,
that they do not know whether their synthetic spikes are realistic enough for the resulting
scores to mean what everyone treats them as meaning. This project takes that sentence
literally and tests it.

What follows is not a summary of the project. It is the conceptual territory the project
crosses. By the end you should be able to read the Claim Sheet --- the contract, and the
densest document in the repository --- without stopping to look anything up, and to follow
Phase 2 as a participant rather than a spectator. Where a concept connects to a specific
commitment in that contract, I say so at the point the concept arrives, rather than
collecting the connections at the end.

Section 2 builds the domain: what these recordings are, why sorting them is hard, and why
the field ended up in a position where it has to manufacture its own answer key. Section 3
takes the methods one at a time --- hybrid generation, the three realism axes this project
manipulates, and the gate that decides whether a manipulation counted. Section 4 covers
evaluation: how accuracy is measured here, and the statistical spine that turns a pile of
numbers into a decision. Section 5 assembles the pieces into a system and names the
assumptions holding it up, because those are the places where the project could be wrong.

\section{Domain Background}

\subsection{What the probe actually hears}

A neuron fires by moving charged particles across its membrane, and that movement makes a
small, brief voltage disturbance in the salty fluid around it. Put a metal contact nearby
and you can record that disturbance without touching the cell. This is \emph{extracellular}
recording, and its great virtue is that it is survivable: the neuron goes on living and
firing, and you can listen for hours.

Its great vice is that the fluid does not sort your mail. A contact hears every neuron
close enough to be heard, all at once, summed into one wiggling voltage trace. It also
hears the electrical hum of everything further away, which is the noise floor.

For most of the history of the field this was tolerable because electrodes were few. Then
the probes changed. A Neuropixels probe is a single silicon shank carrying hundreds of
recording sites spaced tens of micrometres apart, dense enough that one neuron is heard on
several sites at once
(\href{https://doi.org/10.1038/nature24636}{Jun et al., 2017, \emph{Nature}} --- worth
skimming for the images alone; it is the engineering paper that created the data regime
everything downstream now assumes).

That density is the good news and part of the reason the problem got hard. Because a neuron
is heard on several neighbouring sites at different strengths, the \emph{pattern} across sites
is a spatial fingerprint. It reflects where the cell sits relative to the probe, but also the
cell's anatomy, the surrounding tissue, and the probe geometry. That fingerprint, averaged
over many of a neuron's spikes, is called its \textbf{template}. Templates are the central
object of this project: they are what gets injected, and their realism is what Tier A of the
experiment manipulates.

\subsection{Spike sorting, and why it is genuinely hard}

\emph{Spike sorting} is the step that turns the summed voltage trace into a list of the form
``neuron 17 fired at 12.483 seconds.'' Nearly everything downstream in systems neuroscience
--- what a brain area encodes, how cells coordinate, what changes with learning --- is
computed on top of that list. If the list is wrong, so is everything computed from it, and
usually silently.

Four things make it hard, and each one shows up later as a design constraint.

\textbf{Overlap.} Two nearby neurons can fire within a millisecond of each other, and their
voltage signatures add together into a shape belonging to neither. These are
\emph{collisions}. How well an algorithm untangles them is not a minor implementation
detail --- it is the property on which whole families of sorters visibly diverge
(\href{https://doi.org/10.1523/ENEURO.0105-22.2022}{Garcia, Buccino \& Yger, 2022,
\emph{eNeuro}} --- measures this directly, and reports that template-matching sorters
recover collided spikes far better than density-based ones, with reconstruction errors above
50\% at short lags for the latter). Hold onto this: it is the single most important reason
the project insists on comparing mechanistically different sorters rather than two versions
of one.

\textbf{Drift.} Brains move --- breathing, heartbeat, the slow settling of tissue around an
implant. The probe stays put, so a neuron slides along it, and its fingerprint changes as it
goes. A benchmark that crossed 12 drift scenarios found that even with the best available
correction applied, the number of well-recovered neurons fell from 165 to as low as 107
(\href{https://doi.org/10.1523/ENEURO.0229-23.2023}{Garcia et al., 2024, \emph{eNeuro}} ---
useful for calibrating scale: it tells you how large an effect has to be before it competes
with a problem the field already treats as serious).

\textbf{Variability within one neuron.} A neuron's spikes are not identical copies. In
particular, when a cell fires several spikes in quick succession --- a \emph{burst} --- the
recorded amplitude of each successive spike \emph{shrinks}. This matters enormously later,
so hold it: a real neuron's waveform depends on what that neuron just did.

\textbf{No answer key.} There is no independent record of which neuron actually fired in a
real recording. This is not a gap in the tooling; it is the structure of the situation. And
it is where the project lives.

\subsection{Three ways to manufacture an answer key}

Since nature does not supply ground truth, the field manufactures it, and there are exactly
three strategies. The review that lays them out side by side is
\href{https://doi.org/10.1088/2516-1091/ac6b96}{Buccino, Garcia \& Yger, 2022,
\emph{Progress in Biomedical Engineering}} (the best single orientation to modern spike
sorting if you want one source beyond this guide).

\textbf{Paired recordings} put a second electrode inside a single neuron while the probe
listens outside. The inside electrode is unambiguous, so you know that cell's spike times
exactly. This is the gold standard and it is nearly useless at scale: you typically get
\emph{one} verified neuron per recording, and only from the large, sturdy cells that
tolerate being impaled. Real, but rare and biased.

\textbf{Fully synthetic recordings} simulate everything --- neurons, tissue, noise. You get
every neuron's truth for free. The cost is that if the simulation is unrealistic in some
respect, every unit in the dataset is unrealistic in that same respect, all at once, and
nothing in the data will tell you.

\textbf{Hybrid recordings} are the compromise the field actually grades on, and the subject
here. Take a genuine recording --- with its real noise, real drift, real firing neurons,
real artifacts --- and inject a handful of synthetic neurons on top. The injected ones come
with a known answer key by construction; the surrounding reality is not simulated because it
was never synthetic to begin with.

That compromise is why hybrid data won, and it is also exactly where the vulnerability sits.
Only the injected units are graded, so the grade describes how findable \emph{the injected
units} were. If they are systematically easier to find than real neurons, every sorter
scores too high. Worse --- and this is the real target --- if they are easier for one
\emph{kind} of algorithm than another, then the comparison between algorithms is partly
measuring the generator that made the fake spikes rather than the algorithms.

There is already published reason to think this is not paranoia. A benchmark spanning
roughly 35,000 ground-truth units found that synthetic ground truth produces a
systematically different error signature from real paired ground truth, and its authors
attributed this to simulations not yet reproducing the firing and noise statistics of real
recordings (\href{https://doi.org/10.7554/eLife.55167}{Magland et al., 2020, \emph{eLife}}
--- also the source of the accuracy metric used in Section 4, so it is worth a look). The
same worry appears in a 2018 toolbox paper, whose authors wrote plainly that it was not
clear whether their simulated data reproduced the conditions of actual recordings
(\href{https://doi.org/10.7554/eLife.34518}{Yger et al., 2018, \emph{eLife}}).

So the concern is roughly a decade old and has been stated repeatedly by the people building
the tools. What nobody has done is the controlled version: hold everything else fixed, change
exactly one realism property, and measure what moves. That sentence is this project.

\subsection{The specific sentence being tested}

The anchor is \href{https://doi.org/10.7554/eLife.110170.3}{Buccino et al., 2026,
\emph{eLife}} --- the paper describing the pipelines this project puts under test, and the
one whose configuration the control arm reproduces. In its Limitations section its authors
wrote:

\begin{quote}
``It remains to be tested whether generating more realistic hybrid recordings will have any
effect on spike sorting accuracy.''
\end{quote}

And immediately before it, naming a specific deficit:

\begin{quote}
``although we used different templates for each probe type, we did not choose templates to
match the brain region of the original recording.''
\end{quote}

Two things about this are worth internalising, because they shape the project's posture.

First, this is a limitation the authors identified in their own work, voluntarily, in print.
The project is therefore not a critique. It is taking up an invitation. The Claim Sheet
commits to writing the result as ``we tested the limitation you named, and here is what we
found,'' never as ``you were wrong'' --- and that holds regardless of which way the result
comes out.

Second, \textbf{both possible answers are worth publishing}, which is a genuinely unusual
and comfortable position. If realism does not move the comparison between sorters, the
field's default answer key is validated and a stated open question closes. If it does move
it, then published sorter comparisons built on hybrid data have been partly grading the
generator, and every benchmark that used one inherits that. There is no result here that is
a disappointment --- only a result that is too imprecise to be either, which Section 4
addresses head-on.

\section{Core Methods}

\subsection{How a hybrid recording is actually built}

To manipulate realism you need to know where the seams are. Building one hybrid recording
takes four steps.

\textbf{Choose donor templates.} You need spatial fingerprints to inject. Rather than invent
them, the pipeline draws from a published library of templates extracted from real
recordings. Each carries metadata: its amplitude, its signal-to-noise ratio, its depth along
the probe, which dataset it came from --- and, crucially for this project, a
\texttt{brain\_area} label.

\textbf{Choose spike times.} When will each injected neuron fire? The standard pipeline draws
times from a \emph{Poisson process}: a coin flipped very fast, with a fixed average rate and
no memory. Each moment is independent of every other. It is the mathematically simplest way
to generate events in time, and --- as Section 3.3 argues --- it is where most of the
missing realism lives.

\textbf{Place and rescale.} The template is positioned at a chosen depth and scaled so its
amplitude lands in a target range (the anchor used 50--200 $\mu$V), because an injected unit
that is far louder or fainter than the surrounding real cells would be trivially findable or
invisible.

\textbf{Add.} The scaled template is summed into the real voltage trace at each chosen time.
The recording is now a hybrid: real everywhere, plus a handful of neurons whose every spike
time you wrote down.

A useful contrast makes the design choice visible. A different tool, SHYBRID, builds hybrid
data by \emph{relocating a real unit} --- reusing an observed unit's spike times at a new
position on the probe
(\href{https://doi.org/10.1007/s12021-020-09474-8}{Wouters, Kloosterman \& Bertrand, 2021,
\emph{Neuroinformatics}}). Its
\href{https://github.com/jwouters91/shybrid/blob/master/hybridizer/threads.py}{source code}
reuses the observed spike times after a fixed shift and each spike's fitted template amplitude;
the new insertion train separately receives random sub-sample jitter by default. That preserves
more observed train and amplitude structure than an average-template-plus-Poisson construction;
it does not imply that every per-spike waveform shape is transported unchanged. \emph{The
field has already made materially different choices about which structure to preserve},
which is a decent sign that the question is worth an experiment.

\subsection{Choosing what to make more realistic}

``More realistic'' is not a single dial, so the project picks three specific properties,
tests them one at a time, and never varies two at once --- because if you move two things and
the score changes, you have learned nothing about which one did it. In the Claim Sheet these
are Tiers A, B, and C, and they are ordered by how much has to be built before they can be
measured.

\textbf{Tier A --- region-matched templates.} Injected templates currently come from whatever
brain region the library happens to offer, not the region the host recording came from. Is
that a real mismatch or a cosmetic one? It is real: waveform shapes differ enough between
brain regions that a neuron's region can be partly \emph{classified} from its multi-channel
waveform (\href{https://doi.org/10.1152/jn.00680.2018}{Jia et al., 2019, \emph{Journal of
Neurophysiology}} --- the empirical warrant for this axis; it also introduces the
regular-spiking/fast-spiking distinction that shows up whenever waveform shape is discussed).
Because the library already carries a region label, fixing this requires \emph{no new
generator mechanism} --- it is a constrained selection and balance problem over existing
rows. Tier A is measurement rather than generator engineering, and it is first because it is
cheapest.

\textbf{Tier B --- population-coupled firing.} Real neurons do not fire independently of
their surroundings; population activity rises and falls together, and individual rates ride
those waves. Poisson spike times are, by construction, deaf to this. The fix is to estimate
how busy the local population is over time and make injected rates follow it. This is the
anchor authors' \emph{own} proposed remedy, and there is prior art --- Kilosort4's hybrid
benchmark already modulates intervals by local population rate in 100 ms bins. Modest new
code.

\textbf{Tier C --- bursting with history-dependent waveform shrinkage.} Recall from Section
2.2 that a real neuron's spikes shrink across a burst. A fixed average template presents
precisely the same shape no matter what the neuron just did --- which is not a small
idealisation but the removal of a well-documented mechanism. In hippocampal CA1 pyramidal
cells, complex-spike bursts occur at intervals of 6 ms or less with amplitude decreasing
across the burst, and bursts are suppressed following recent activity
(\href{https://doi.org/10.1016/S0896-6273(01)00447-0}{Harris et al., 2001, \emph{Neuron}} ---
supplies the actual numbers Tier C has to hit). That this matters for \emph{sorting}
specifically is not a guess: a sorter built two decades ago improved its results by
explicitly modelling burst-dependent amplitude attenuation
(\href{https://doi.org/10.1152/jn.00227.2003}{Pouzat et al., 2004, \emph{Journal of
Neurophysiology}}). This mechanism is genuinely absent from the pipeline under test, and it
has to be built before it can be measured.

Two honesties are written into the contract about this list, and they are worth knowing
because they are the kind of thing a project can quietly slide away from.

The first is the \textbf{measurement/engineering split}. Tier A is measurement. Tiers B and C
are engineering in the service of measurement. That is still research, but the Claim Sheet
says so plainly rather than presenting the whole thing as pure measurement.

The second is sharper. \textbf{The project's prior is that Tier A is most likely to move
absolute scores and least likely to move the comparison between sorters.} Changing static
waveform shape changes what every sorter's front end sees, plausibly in broadly similar ways.
The temporal tiers hit collision handling --- where, per Section 2.2, sorter families
actually diverge. Since Tier A runs first and is cheapest, there is a live risk of appearing
to have answered the headline question when the project has answered its easiest version.
The contract therefore states that \emph{the project may not conclude on Tier A alone}. That
commitment exists to stop a convenient early stopping point from becoming the result.

\subsection{The manipulation check: the gate before the experiment}

Here is a failure mode that would quietly ruin the whole project.

Suppose Tier C is implemented, run, and the sorters' scores do not move. Two explanations
fit: realism genuinely does not matter, or the bursts that were injected were feeble,
biologically implausible, or subtly broken. Those two conclusions are opposite in meaning and
identical in appearance. Nothing measured \emph{after} the sorters run can separate them.

So the project inserts a gate. Before any sorter time is spent on a tier, the generated data
must be inspected and shown to carry the intended property at a realistic magnitude ---
the injected inter-spike intervals really do reach the biological range, the amplitude really
does attenuate with recent firing, and the quantities that were supposed to stay fixed really
did stay fixed. Only then does sorting begin. \textbf{If the check fails, no sorter runs at
all}; the failure is diagnosed and published as a finding about the generator.

This is the most important structural idea in the design, and it is worth seeing why it is
not merely good hygiene. It is what makes a \emph{negative} result mean anything. Without the
gate, ``realism did not matter'' is indistinguishable from ``we implemented realism badly'' ---
and since a negative result is the likelier and arguably more useful outcome here, the gate is
what keeps the project's most probable finding from being worthless.

The same logic drives the \textbf{covariate balance} requirement. If the region-matched
templates happened to be louder on average than the region-unaware ones, then any measured
effect could be an amplitude effect wearing a region label. So amplitude, signal-to-noise,
probe geometry, placement, and even how many different source datasets each arm's templates
came from are all balanced across arms and checked. In the contract this is called a gate
rather than a hope, which is the right register: it is verified, not assumed.

\subsection{Why the comparison is paired}

The anchor study estimated roughly 870 and 846 single-workstation hours for its two sorter
comparisons. This project runs on one shared desktop that other projects are using at the
same time. It cannot buy confidence with volume.

It buys it with \emph{pairing} instead. If you want to know whether a change matters, the
statistically expensive way is to run a large group with and a large group without and
compare averages, absorbing all the variation between individuals. The cheap way is to test
the \emph{same} individuals under both conditions and look at each one's change. Most of the
noise --- the fact that some injected units are intrinsically easier to find than others ---
cancels, because it is present identically on both sides of the subtraction.

Concretely: the same host recording, the same placements, the same amplitude targets, the same
unit count, the same total spike counts, with exactly one property changed. For Tiers B and C
it is literally the same donor template in both arms.

Tier A cannot do this, and the reason is instructive. Changing donor \emph{region}
necessarily changes donor \emph{identity} --- a region-matched template and a region-unaware
one are different templates by definition. So Tier A pairs \emph{matched slots} rather than
identical units: two donors deliberately selected to resemble each other on every measured
nuisance property, injected with the same spike train at the same place with the same
rescaled amplitude. The pairing claim is about the matched slot, not about a fictitiously
unchanged unit. This distinction was caught during review of the Claim Sheet, and it is a
good example of what the review cycle between the two agents is for.

\section{Evaluation Approach}

\subsection{Scoring one sorter on one recording}

After a sorter runs, its output is compared against the known injected spike times. For each
injected unit, three counts fall out: \emph{matches} (found correctly), \emph{misses} (real
spikes the sorter never reported), and \emph{false positives} (spikes the sorter reported
that were not there). The field's standard accuracy metric combines all three:

\[
\text{accuracy} \;=\; \frac{n_{\text{match}}}{n_{\text{match}} + n_{\text{miss}} + n_{\text{fp}}}
\]

In words: of everything that either happened or was claimed to happen, what fraction did the
sorter get right? It ranges from 0 to 1, and it is deliberately unforgiving --- inventing
spikes is penalised just as missing them is. The convention comes from
\href{https://doi.org/10.7554/eLife.55167}{Magland et al., 2020}, which also established the
practice of reporting the \emph{count of units scoring above 0.8} alongside the average.

That second convention matters more than it looks. An average can hide the difference between
a manipulation that nudged every unit slightly and one that pushed a few units across a line
that a researcher actually acts on. This project reports both, and it also breaks results out
by signal-to-noise ratio --- because an effect concentrated in the faint, hard units could
leave the average almost untouched and still change every practical conclusion.

\subsection{The quantity the project actually cares about}

Now the conceptual centre of the whole design, and the one place where notation earns its
keep.

Take one sorter. For each paired unit, subtract its accuracy in the control arm from its
accuracy in the realistic arm. Call that sorter's \textbf{realism effect}:

\[
\Delta_s \;=\; A(\text{realistic}) - A(\text{control})
\]

where $A$ is accuracy and $s$ names the sorter. If $\Delta_s$ is negative, that sorter did
worse when the injected spikes were made more realistic --- which is the expected direction,
since realism should make the planted spikes harder to find.

But $\Delta_s$ alone is the \emph{absolute-score} question. If every sorter drops by the same
amount, published accuracy numbers are too optimistic --- yet every ranking, every ``sorter A
beat sorter B,'' survives untouched. Since hybrid benchmarks are used mostly to \emph{choose
between} sorters, that would be a real but limited finding.

The question that threatens the field's actual practice is whether realism affects sorters
\emph{differently}. So the primary quantity is the difference between two sorters' realism
effects --- a \textbf{difference in differences}, also called an \textbf{interaction}:

\[
I(s_1, s_2) \;=\; \Delta_{s_1} - \Delta_{s_2}
\]

Read it aloud: \emph{how much more did sorter 1's accuracy change under realism than sorter
2's?} A negative value means realism hurt sorter 1 more; a positive value means it hurt
sorter 2 more (or helped sorter 1 more). If $I$ is near zero, realism moved both sorters
together and the comparison between them is stable --- though their absolute scores may
still be realism-sensitive. If $|I|$ is large, the gap between the two sorters depends on how
realistic the fake spikes were, which means the benchmark has been partly grading the
generator.

That is the project's target. Both the individual $\Delta_s$ values and the interaction $I$
are reported separately, because they answer different questions and a project that reported
only one of them would be hiding half of what it found.

\subsection{How much is enough to matter}

A number without a threshold is not yet a decision. Two thresholds are fixed \emph{before any
result is seen}, which is the entire point --- a threshold chosen after seeing the data is not
a threshold, it is a rationalisation.

For \textbf{absolute scores}, a change of 0.05 accuracy (five points) is called material. The
reasoning is concrete rather than conventional: the field counts units above 0.8, so a shift
of that size near the threshold moves units across it and changes a reported number.

For the \textbf{comparison}, the margin is measured against the sorter gap actually observed
in this experiment. Let $G_0$ be the average gap between the two sorters in the control arm.
The margin is

\[
T \;=\; \max(0.05,\; 0.5 \times |G_0|)
\]

Two ideas are packed in there. The half-gap term asks whether realism moved a meaningful
\emph{fraction} of the comparison that actually exists; the fraction is a declared judgement,
not a theorem that every half-gap shift changes a reader's decision. The five-point floor
prevents an absurdity: if two sorters were nearly tied in the control arm,
$0.5 \times |G_0|$ would be almost zero, and a trivial wobble would qualify as
decision-relevant.

\subsection{Confidence, and the honest third answer}

Every number here is an estimate from a finite experiment, so each needs an interval around it
rather than a bare point. The tool is the \textbf{bootstrap}: rather than assume a formula for
how much a measurement would wobble if the experiment were repeated, you simulate the
repetition by resampling your own data many times and watching how much the answer moves
(\href{https://doi.org/10.1214/aos/1176344552}{Efron, 1979, \emph{Annals of Statistics}} ---
the original; its appeal is that it replaces a distributional assumption with computation).

The detail that matters is \emph{what gets resampled}. Ten injected units inside a single run
are not ten independent experiments --- they share a recording, a noise environment, and a
randomization block. Treating them as ten independent repetitions would produce an interval
far too narrow and a confidence far too high. The bootstrap therefore mirrors the design:
paired unit or donor-slot observations are resampled within randomization blocks while their
arm and sorter pairing stays intact; blocks are resampled within hosts; and, once there is
more than one host, hosts are resampled at the top level and weighted equally. Repeated donor
identities stay in one resampling cluster. With one host, the interval is explicitly
conditional on that host and cannot manufacture cross-host generalization. Every resample
recomputes $G_0$, $T$, and the interaction rather than treating the threshold as known
exactly.

The Claim Sheet records the comparative decision with one final quantity,
$D = |I| - T$: the interaction's magnitude minus the materiality margin. A 95\% interval for
$D$ wholly above zero, together with an interaction interval excluding zero, is a bounded
positive. An interval for $D$ wholly below zero is a bounded negative. Crossing zero is
inconclusive. The absolute value makes a near-zero interaction resample slightly upward on
average, so the rule is deliberately conservative about declaring a negative: it will call a
real null inconclusive before it calls noise a settled null.

Which produces three possible verdicts, not two:

\begin{itemize}
\item \textbf{Bounded positive} --- the interaction's magnitude is confidently larger than the margin.
Realism changes the comparison between sorters.
\item \textbf{Bounded negative} --- the interaction's magnitude is confidently smaller than the margin.
Realism does not change the comparison, within the tested boundaries. \emph{This is a success,
published on the same terms as a positive.}
\item \textbf{Inconclusive} --- the interval is too wide to say which. The honest report is
``not resolved at this precision,'' with the achieved bound stated.
\end{itemize}

The third one deserves emphasis, because it is where a project like this is most likely to
mislead people without meaning to. A wide interval that happens to be centred near zero
\emph{looks} like ``realism does not matter.'' It is not. It is ``we could not tell.'' The
Claim Sheet states in advance that a null with a wide interval will never be reported as
evidence that realism does not matter --- a commitment made now precisely because it would be
much harder to honour later, when a tidy conclusion is within reach and the compute budget is
spent.

\subsection{The negative control}

One more guard, and it answers a question a sceptical reader would ask immediately: how much
apparent interaction would this design produce \emph{if nothing were manipulated at all}?

The answer is measured rather than argued. Pairs of arms are generated under the \emph{same}
nominal condition --- same settings, different random draws of which templates were selected,
which spike times were rolled, where things landed --- and put through the identical
comparison. Any interaction they show is pure noise from selection and seeds, because no
realism property differed between them. That spread becomes a reference band drawn behind the
real result.

Its role is precise, and reviewing the Claim Sheet tightened it. The band is a
\emph{diagnostic} --- it shows whether the machinery can manufacture an effect of the observed
size out of nothing --- and it is emphatically not a second significance test. The decision is
made by the interval and the pre-declared margin. A shaded band that quietly becomes a truth
threshold is a way of testing the same data twice, and the contract now says so explicitly.

These pseudo-arms have a cost that is easy to overlook: they must be \emph{sorted}, not merely
generated, which doubles the sorter time the primary run needs. That is now written into the
compute budget rather than discovered mid-experiment.

\section{How It All Fits Together}

\subsection{The pipeline as a system}

Read end to end, the machine has six stages, each handing something specific to the next.

\textbf{A real recording} enters, supplying genuine noise, genuine drift, and genuinely
firing neurons. \textbf{Donor templates} are selected from the library --- region-matched or
region-unaware, balanced on every nuisance property. \textbf{Spike times} are generated ---
Poisson for the control, population-coupled or bursty for the realism arms. \textbf{Injection}
sums scaled templates into the trace, producing a hybrid recording with a known answer key.
\textbf{Sorters} run on it, blind to which arm they are in, on pinned default parameters that
are never tuned per condition --- because tuning on a condition leaks the condition into the
result. \textbf{Comparison} scores the output against the answer key, and the scores flow into
the paired differences, the interaction, and the bootstrap interval.

Two of those stages are gates rather than steps. The manipulation check sits between generation
and sorting, and can stop everything. The covariate balance check sits inside donor selection
and can reject a host recording entirely.

\subsection{Where the leverage is}

If you want to know which choices actually determine the outcome, they are these.

\textbf{The sorter panel.} Kilosort4 plus at least one mechanistically \emph{different} CPU
sorter. Comparing two versions of Kilosort would share a family and a drift-correction lineage
--- the pairing least likely to show a difference, which would make a null result look
designed-for. Since collision handling is where families diverge, the panel has to span
families for the temporal tiers to have anything to reveal.

\textbf{The host recording and where on the probe the injections go.} A probe passes through
several brain areas, so the project cannot label a whole recording with one region; it pins an
anatomical injection zone and matches donors to the local label there. The contract also
makes explicit that the same host should be used across all three tiers, because its
interpretation uses cross-tier reasoning, and a host that changes between tiers turns every
such statement into a comparison across recordings as well as across axes.

\textbf{Compute, which is the binding constraint.} The machine is shared with other Dandelion
projects that are not coordinated with this one. Free memory is a measurement that changes
underneath the work rather than a budget allocated to it --- at two separate moments during
early sessions the machine had under 4.5 GiB free, against a known-good Kilosort4 run that
peaked at 29.3 GiB. Started at either moment, that run would have failed slowly and
confusingly. Hence the rule that every heavy step measures free memory immediately beforehand
and does not start what does not fit.

\subsection{The load-bearing assumptions}

These are the places where the project could be wrong in ways that would matter, listed so
you can hold it to them.

\textbf{That the manipulation check can actually be passed.} Everything downstream assumes
that bursting and population coupling can be implemented at genuinely biological magnitudes.
If they cannot, the project reports a generator failure --- which is a real finding, but not
the one it set out for.

\textbf{That the biological numbers transfer.} Tier C's burst parameters come from
hippocampal CA1 pyramidal cells. They are not generic brain-wide biology. Using them elsewhere
without fresh evidence would make Tier C a synthetic stress test rather than a
biological-realism test, and the contract requires it to be labelled as such if that happens.
This constraint is what makes CA1 the leading \emph{candidate} host zone --- it is the one
place where Tier C's biological prior and Tier A's donor-availability screen currently
overlap. It is not yet a selection: its worst-case donor exclusion leaves only six templates,
so the exact host source, placement, and covariate-balance checks still decide whether it is
viable.

\textbf{That one host does not generalise.} A result from one recording, one injection zone,
one probe type is a result about that configuration. Widening is a rung on the ladder, and
absent it the claim carries the boundary explicitly.

\textbf{That Tier B cannot fully attribute its own result.} Kilosort4's own benchmark already
uses population-rate coupling. Adding it to the SpikeInterface pipeline therefore moves the
test data toward Kilosort4's home turf. If Kilosort4 gains under Tier B, ``more robust to
realistic firing'' and ``developed against data that already had this property'' cannot be
separated by this design. That is declared in advance as inconclusive-on-attribution rather
than discovered afterwards and argued about.

\subsection{What to watch for as Phase 2 runs}

Four things will tell you whether the project is on its rails.

Does the \textbf{manipulation check} pass, and at what magnitude? If the intended property is
present but below the pre-declared biological range, the biological-realism gate has not
passed; any downstream stress-test result is bounded to the magnitude actually achieved.

Does the \textbf{pilot} admit a second and third sorter within its declared budget? If only two
fit, the panel narrows and the project has pre-committed to reporting that as a limitation
rather than quietly proceeding.

Is the \textbf{interval} narrow enough to decide anything? This is where the compute constraint
becomes a scientific constraint, and where ``inconclusive'' is a live and honourable outcome.

Does the \textbf{negative control band} stay small? If shuffling seeds produces interactions as
large as the realism manipulation does, the design cannot resolve the question at this scale,
and the honest move is more replication rather than a conclusion.

\subsection{Reading the Claim Sheet from here}

You now have what the contract assumes. Its fifteen slots run from what the project works with
through to how it could be wrong. If you read only part of it, read Slots 11, 12, and 13 --- the
pre-declared shapes of success, failure, and inconclusiveness. Everything in Sections 4.3 and
4.4 above is the vocabulary those three slots are written in, and they are the reason the rest
of the document had to be written before the work started rather than after it.

\enlargethispage{2\baselineskip}
One closing note. Almost every safeguard here --- the gate, the negative control, pre-declared
thresholds, the refusal to conclude on the cheapest axis, and the wide-interval outcome ---
makes tidy answers harder. Deliberately. The question is worth asking only if its answer can
be trusted; studies like this usually fail \emph{convincingly}, not obviously.

\end{document}
